\section{Exact full-face modes and the global correction}
\label{sec:terminal-modal-exact}

Let $r_k^x=r(p_k)$ and $r_k^v=r(v_k)$ after entry to the full face, with entry
indexed by $k=0$.  Whenever the projection in
\cref{eq:terminal-modal-step} is inactive, the sibling note proves
\begin{align}
 r_{k+1}^x&=\frac{B}{1+q}(r_k^x+q r_k^v),
 &q r_{k+1}^v&=r_{k+1}^x-(1-q)r_k^x,
 \label{eq:terminal-modal-two-state}\\
 r_{k+1}^x&=B\left(\frac2{1+q}r_k^x
 -\frac{1-q}{1+q}r_{k-1}^x\right),
 &B&=I-\widehat Q.
 \label{eq:terminal-modal-second-order}
\end{align}

For $0\le h\le m$, define
\begin{equation}
 u_h(j)=\cos\frac{h\pi j}{m},\qquad 0\le j\le m.
 \label{eq:terminal-modal-cosines}
\end{equation}
Use the reversible inner product
$\langle z,w\rangle_D=\sum_{j=0}^m d_jz_jw_j$.

\begin{lemma}[Weighted path cosines; proved here]
\label{lem:terminal-modal-cosines}
The vectors $u_h$ are an orthogonal eigenbasis of $P_{\rm row}$ with
eigenvalues $\mu_h=\cos(h\pi/m)$.  Their squared norms are
\begin{equation}
 N_h:=\norm{u_h}_{D}^{2}
 =\begin{cases}
 2m,&h\in\{0,m\},\\
 m,&1\le h\le m-1.
 \end{cases}
 \label{eq:terminal-modal-norms}
\end{equation}
Consequently every vector $z$ has coefficients
$\widehat z_h=\langle z,u_h\rangle_D/N_h$, and
\begin{equation}
 \norm{z-\widehat z_0\one}_D^2
 =m\sum_{h=1}^{m-1}\widehat z_h^2+2m\widehat z_m^2.
 \label{eq:terminal-modal-parseval}
\end{equation}
\end{lemma}

\begin{proof}
The endpoint equations are
$(P_{\rm row}u_h)_0=u_h(1)=\mu_hu_h(0)$ and similarly at $m$.
At an interior point, the cosine addition formula gives
$(u_h(j-1)+u_h(j+1))/2=\mu_hu_h(j)$.  Reversibility,
$DP_{\rm row}=A=A^\top$, makes distinct modes $D$-orthogonal.  Finally,
the standard discrete-cosine sums with endpoint weights one and interior
weights two give \eqref{eq:terminal-modal-norms}; Parseval follows.
\end{proof}

Let $a_{h,k}=\langle r_k^x,u_h\rangle_D/N_h$ and set
\begin{equation}
 \phi_h=\frac{h\pi}{2m},\qquad
 \varrho_h=(1-q)\cos\phi_h.
 \label{eq:terminal-modal-phase-radius}
\end{equation}

\begin{lemma}[Literal modal solution; proved here]
\label{lem:terminal-modal-solution}
For $1\le h\le m-1$, every projection-inactive full-face trajectory obeys
\begin{equation}
 a_{h,k}=\varrho_h^k
 \left(C_h\cos(k\phi_h)+D_h\sin(k\phi_h)\right),
 \label{eq:terminal-modal-solution}
\end{equation}
where
\begin{equation}
 C_h=a_{h,0},\qquad
 D_h=\frac{a_{h,1}/\varrho_h-C_h\cos\phi_h}{\sin\phi_h}.
 \label{eq:terminal-modal-quadratures}
\end{equation}
The constant mode has a repeated root $1-q$ and is irrelevant to residual
range.  The $h=m$ mode is killed after one application of $B$.
\end{lemma}

\begin{proof}
On mode $h$, the eigenvalue of $B$ is
$(1-q^2)\cos^2\phi_h$.  Thus
\cref{eq:terminal-modal-second-order} has characteristic polynomial
\[
 z^2-2(1-q)\cos^2\phi_h\,z
 +(1-q)^2\cos^2\phi_h,
\]
whose roots are $\varrho_he^{\pm i\phi_h}$.  Matching $k=0,1$ gives
\eqref{eq:terminal-modal-quadratures}.  The two endpoint cases follow by
direct substitution.
\end{proof}

The apparent sine loss in \eqref{eq:terminal-modal-quadratures} can be
removed exactly at entry.  Define the entry centered-residual coefficient
\begin{equation}
 V_h:=\frac{\langle r_0^v,u_h\rangle_D}{N_h}.
 \label{eq:terminal-modal-entry-velocity}
\end{equation}

\begin{lemma}[Exact entry-velocity quadrature; proved here]
\label{lem:terminal-modal-entry-velocity}
If the first full-face projection is inactive, then, for
$1\le h\le m-1$,
\begin{equation}
 D_h=q\cot(\phi_h)V_h.
 \label{eq:terminal-modal-entry-velocity-quadrature}
\end{equation}
\end{lemma}

\begin{proof}
The eigenvalue of $B$ on $u_h$ is
$\beta_h=(1-q^2)\cos^2\phi_h$.  The first equation in
\eqref{eq:terminal-modal-two-state} therefore gives
\[
 a_{h,1}
 =\frac{\beta_h}{1+q}(C_h+qV_h)
 =(1-q)\cos^2\phi_h(C_h+qV_h).
\]
Because $\varrho_h=(1-q)\cos\phi_h$, substitution in
\eqref{eq:terminal-modal-quadratures} leaves exactly
$qV_h\cos\phi_h/\sin\phi_h$.
\end{proof}

The safe-envelope correction on the full face is
\begin{equation}
 \delta_k=\max\left(0,\frac{\max_j r_{k,j}^x}{\alpha}\right),
 \qquad
 \ell_k=[p_k-\delta_k\one]_+.
 \label{eq:terminal-modal-correction}
\end{equation}
This maximum is global; it is not a seed or frontier surrogate.

\begin{lemma}[Literal range implication; proved here]
\label{lem:terminal-modal-range}
Suppose $p_k-\delta_k\one>0$ coordinatewise.  Put
$R_k=\max_jr_{k,j}^x-\min_jr_{k,j}^x$.  If
\begin{equation}
 R_k>\alpha\tau=\frac{q^3}{5},
 \label{eq:terminal-modal-range-threshold}
\end{equation}
then the one-sided certificate fails at step $k$.
\end{lemma}

\begin{proof}
Because $P_{\rm row}\one=\one$, one has $\widehat Q\one=\alpha\one$.
Unclipping therefore gives
\[
 r(\ell_k)=r_k^x-M_k\one,
 \qquad M_k=\max(0,\max_jr_{k,j}^x).
\]
If $\max_jr_{k,j}^x\ge0$, the minimum corrected residual is exactly $-R_k$.
If the maximum is negative, the correction is zero and
$\min_jr_{k,j}^x< -R_k$.  In either case
\eqref{eq:terminal-modal-range-threshold} leaves an active row below
$-\alpha\tau$.
\end{proof}

Let $\bar r_k=\langle r_k^x,\one\rangle_D/(2m)$ and
$s_k=r_k^x-\bar r_k\one$.  Weighted Popoviciu gives
\begin{equation}
 R_k\ge\frac{2\norm{s_k}_D}{\sqrt{2m}}.
 \label{eq:terminal-modal-popoviciu}
\end{equation}
Indeed, the degree weights divided by $2m$ define a probability distribution,
and the variance of a variable in an interval of length $R_k$ is at most
$R_k^2/4$.  This is why modal $\ell_2$ energy prevents spatial cancellation
without selecting a particular vertex.
